\section{Gaussian operator norms and relative entropy}
\label{sec:operator-entropy}
The entropy lower bound need not pay for every Gaussian parameter. Operator norms give the same output-energy floor while their Gaussian cost depends only linearly on learner width.

For the normalized matrices
\[
 A_1=\sqrt{d/N}\,W^1,\quad A_\ell=W^\ell\ (2\le\ell\le H),
 \quad A_{H+1}=\sqrt N\,v^{\mathsf T},
\]
define
\[
 Q_{\mathrm{op}}(w)=\sum_{\ell=1}^{H+1}\|A_\ell\|_{\mathrm{op}}^2.
\]
The Frobenius energy $Q_N$ is the same sum with Frobenius norms. Under the original independent Gaussian prior, every entry of every normalized block has variance $s=\sigma^2$.

\begin{lemma}[Gaussian operator-norm exponential moment]
\label{lem:operator-mgf}
If $G$ is an $r$ by $c$ matrix with independent standard Gaussian entries, then, for $0<\eta<1$,
\[
 \log\E\exp\{(1-\eta)\|G\|_{\mathrm{op}}^2/2\}
 \le\frac12\log(1/\eta)
       +\frac{1-\eta}{2\eta}(\sqrt r+\sqrt c)^2.
\]
For a Gaussian vector the second term may use its dimension instead of $(\sqrt r+\sqrt c)^2$.
\end{lemma}
\begin{proof}
The operator norm is one-Lipschitz for the Euclidean norm of all matrix entries. Gaussian concentration in exponential-moment form gives, for all real $t$,
\[
 \E e^{t(\|G\|_{\mathrm{op}}-m)}\le e^{t^2/2},
 \qquad m=\E\|G\|_{\mathrm{op}}.
\]
With an independent standard Gaussian $Z$ and $a=1-\eta$, the Gaussian exponential identity and Tonelli yield
\begin{align*}
 \E e^{a\|G\|_{\mathrm{op}}^2/2}
 &=\E_Z\E_G e^{\sqrt aZ\|G\|_{\mathrm{op}}}\\
 &\le\E_Z e^{\sqrt aZm+aZ^2/2}
 =\eta^{-1/2}\exp\{a m^2/(2\eta)\}.
\end{align*}
Finally $m\le\sqrt r+\sqrt c$ by the Gaussian matrix comparison bound, while a Gaussian vector has expected norm at most the square root of its dimension. One derivation of the matrix bound compares the centered processes $u^{\mathsf T}Gv$ and $g^{\mathsf T}u+h^{\mathsf T}v$ on pairs of unit vectors: their increment variances are respectively $2-2\langle u,u'\rangle\langle v,v'\rangle$ and $4-2\langle u,u'\rangle-2\langle v,v'\rangle$, and their difference is $2(1-\langle u,u'\rangle)(1-\langle v,v'\rangle)\ge0$. Gaussian comparison then gives $\E\|G\|_{\mathrm{op}}\le\E\|g\|+\E\|h\|\le\sqrt r+\sqrt c$.
\end{proof}

\begin{lemma}[Linear-width entropy cost]
\label{lem:operator-entropy}
For any probability law $\rho$ absolutely continuous with respect to the Gaussian neural prior $\mu_s$,
\[
 2s\KL(\rho\|\mu_s)
 \ge(1-\eta)\E_\rho Q_{\mathrm{op}}-s\mathcal B_\eta(H,N,d),
\]
where
\[
 \mathcal B_\eta(H,N,d)
 =(H+1)\log(1/\eta)
 +\frac{1-\eta}{\eta}\{(4H-1)N+2(d+1)\}.
\]
In particular $\mathcal B_\eta=O_\eta(H(N+d))$.
\end{lemma}
\begin{proof}
The normalized Gaussian blocks are independent. Apply Lemma~\ref{lem:operator-mgf} to each block and add their logarithmic moment bounds. The first block contributes at most $2(N+d+1)$ to the sum of squared expected norm bounds, the $H-1$ internal blocks contribute at most $4(H-1)N$, and the vector readout contributes at most $N$. Their sum is $(4H-1)N+2(d+1)$. The entropy variational inequality with test function $(1-\eta)Q_{\mathrm{op}}/(2s)$ gives the result, including by truncation if an expectation is infinite.
\end{proof}

\begin{corollary}[A posterior objective bound at every super-mean-field exponent]
\label{cor:operator-posterior-objective}
For the exact Gaussian-likelihood posterior $\Pi$ at parameter $\lambda>0$, let
$F=-2\lambda s\log\E_{\mu_s}\exp[-\widehat R/(2\lambda s)]$.
Then
\[
 \E_\Pi[\widehat R+\lambda(1-\eta)Q_{\mathrm{op}}]
 \le F+\lambda s\mathcal B_\eta(H,N,d).
\]
The output-energy floor is unchanged: for every zero-origin one-Lipschitz activation and every input,
\[
 |z_w(x)|\le\frac{\|(1,x)\|}{\sqrt d}
       \left(\frac{Q_{\mathrm{op}}(w)}{H+1}\right)^{(H+1)/2}.
\]
For the support $\mathcal X_D$ the leading factor is $\sqrt{2/d}\le1$. In particular, if the empirical prediction norm exceeds a fixed $t>0$, then
\[
 Q_{\mathrm{op}}(w)\ge(H+1)t^{2/(H+1)}=H+O_t(1).
\]
For every fixed $\gamma>1$, every fixed $\eta>0$, and a bounded likelihood-temperature range $0<\lambda\le\Lambda$, the objective error $\lambda s\mathcal B_\eta$ can be made arbitrarily small uniformly at all $s\le N^{-\gamma}$ and all $N\ge e^{C_\gamma H}$, with $d=O(H)$.
\end{corollary}
\begin{proof}
The exact Gibbs variational identity is
$F=\E_\Pi\widehat R+2\lambda s\KL(\Pi\|\mu_s)$.
Combine it with Lemma~\ref{lem:operator-entropy}. The deterministic output bound follows by applying the activation's one-Lipschitz, zero-origin property at each layer, using the normalized operator norms, and applying arithmetic--geometric mean to their squares. The last assertion follows from
$s\mathcal B_\eta=O_\eta(H(N+d)N^{-\gamma})$, which tends to zero uniformly over the stated widths after increasing $C_\gamma$ and the depth threshold.
\end{proof}
